% BEGIN S2
\par\medskip\noindent\textbf{2 Diameter-3 for Unit Balls in 3D}\par
% BEGIN Thmtheorem9
\par\medskip\noindent\textbf{Theorem 9.}\par
Assuming the OV hypothesis, there is no $O(n^{2-\varepsilon})$ time algorithm for deciding if the intersection graph of a given set of $n$ unit balls in $\mathbb{R}^{3}$ has diameter at most 3.
\par
\par
\par
% END Thmtheorem9
% BEGIN S2.p1
The reduction is from Diameter-2 for $n$-vertex sparse tripartite graphs $G=(A\cup B\cup C,E)$, where the number of edges $m$ is $\tilde{O}(n)$. It is well-known that solving Diameter-2 for sparse tripartite graphs in truly-subquadratic time would break the OV hypothesis [19].
\par
\par
% END S2.p1
% BEGIN Thmtheorem10
\par\medskip\noindent\textbf{Lemma 10 (Diameter-2 sparse tripartite graphs).}\par
There is no $O(n^{2-\varepsilon})$ time algorithm for deciding if a sparse tripartite graph has diameter at most $2$, assuming the OV hypothesis.
\par
\par
\par
% END Thmtheorem10
% BEGIN S2.p2
In this section, we work with balls of radius $1/2$ instead of $1$. Let $\varepsilon=\Theta(1/n^{4})$ and $\delta=\Theta(1/n^{3})$. First, we map every vertex $v$ in $A\cup C$ to a distinct number, also denoted by $v$, in $[\varepsilon,2\varepsilon]$, such that $\min_{a_{1},a_{2}\in A:\,a_{1}\neq a_{2}}|a_{1}-a_{2}|\geq\Omega(\varepsilon/n)$ and $\min_{c_{1},c_{2}\in C:\,c_{1}\neq c_{2}}|c_{1}-c_{2}|\geq\Omega(\varepsilon/n)$. Vertices of $B$ are mapped to distinct numbers in $[0.9,1]$ such that $\min_{b_{1},b_{2}\in B:\,b_{1}\neq b_{2}}|b_{1}-b_{2}|\geq\Omega(1/n)$. Then we create four sets of balls as follows.
\par
\par
% END S2.p2
% BEGIN S2.p3
\begin{enumerate}
\item[]  1. 
For every vertex $a\in A$, we add a ball centered at $s_{a}=(1+a\cos\delta,a\sin\delta,0)$.
\par
\par
\item[]  2. 
For every edge $(a,b)\in(A\times B)\cap E$, we add a ball centered at $p_{ab}=(1+a\cos\delta-b\sin\delta,a\sin\delta+b\cos\delta,\sqrt{1-b^{2}})$.
\par
\par
\item[]  3. 
For every edge $(b,c)\in(B\times C)\cap E$, we add a ball centered at $q_{bc}=(-c,b,\sqrt{1-b^{2}})$.
\par
\par
\item[]  4. 
For every vertex $c\in C$, we add a ball centered at $t_{c}=(-c,0,0)$.
\par
\par
\end{enumerate}
Let $\mathcal{B}_{i}$ be the set of balls created in step $i$ for $i\in[4]$ and $\mathcal{B}\coloneqq\cup_{i}\mathcal{B}_{i}$. Observe that:
\par
\par
% END S2.p3
% BEGIN Thmtheorem11
\par\medskip\noindent\textbf{Observation 11.}\par
Let $\llparenthesis p\rrparenthesis$ be the ball of radius $1/2$ centered at a point $p$. We have:
\par
\begin{enumerate}
\item[]  1. 
$\llparenthesis s_{a}\rrparenthesis\cap\llparenthesis p_{\tilde{a}b}%
\rrparenthesis\not=\emptyset$ if and only if $a=\tilde{a}$.
\par
\par
\item[]  2. 
$\llparenthesis p_{ab}\rrparenthesis\cap\llparenthesis q_{\tilde{b}c}%
\rrparenthesis\not=\emptyset$ if and only if $b=\tilde{b}$.
\par
\par
\item[]  3. 
$\llparenthesis p_{bc}\rrparenthesis\cap\llparenthesis t_{\tilde{c}}%
\rrparenthesis\not=\emptyset$ if and only if $c=\tilde{c}$.
\par
\par
\item[]  4. 
Every two balls in $\mathcal{B}_{i}$ intersect for every $i\in[4]$.
\par
\par
\item[]  5. 
For any two balls $b_{1}\in\mathcal{B}_{i}$ and $b_{2}\in\mathcal{B}_{j}$ such that $|i-j|\geq 2$, $b_{1}\cap b_{2}=\emptyset$.
\par
\par
\end{enumerate}
\par
\par
% END Thmtheorem11
% BEGIN S2.SS0.SSS0.P0.SPx1
\par\medskip\noindent\textbf{Proof.}\par
% BEGIN S2.SS0.SSS0.P0.SPx1.p1
For item 1, observe that
\par
\[\begin{aligned}
\begin{split}\|s_{a}-p_{\tilde{a}b}\|^{2}_{2}&=((a-\tilde{a})\cos\delta+b\sin%
\delta)^{2}+((a-\tilde{a})\sin\delta-b\cos\delta)^{2}+1-b^{2}\\
&=(a-\tilde{a})^{2}(\cos^{2}\delta+\sin^{2}\delta)+b^{2}(\sin^{2}\delta+\cos^{%
2}\delta)+1-b^{2}\nobreak\ =\nobreak\ (a-\tilde{a})^{2}+1,\end{split}
\end{aligned}\]
which is exactly 1 if $a=\tilde{a}$, and at least $1+\Omega(1/n^{10})$ if $|a-\tilde{a}|\geq\Omega(\varepsilon/n)$.
\par
\par
% END S2.SS0.SSS0.P0.SPx1.p1
% BEGIN S2.SS0.SSS0.P0.SPx1.p2
For item 2, observe that
\par
\[\begin{aligned}
\begin{split}&\|p_{ab}-q_{\tilde{b}c}\|^{2}_{2}\\
=&(1+a\cos\delta-b\sin\delta+c)^{2}+(a\sin\delta+b\cos\delta-\tilde{b})^{2}+(%
\sqrt{1-b^{2}}-\sqrt{1-\tilde{b}^{2}})^{2}\\
=&(1+\Theta(\varepsilon)-\Theta(\delta)+\Theta(\varepsilon))^{2}+(\Theta(%
\varepsilon\delta)+b-\Theta(\delta^{2})-\tilde{b})^{2}+(\sqrt{1-b^{2}}-\sqrt{1%
-\tilde{b}^{2}})^{2}\\
=&(1-\Theta(1/n^{3}))^{2}+(|b-\tilde{b}|\pm\Theta(1/n^{6}))^{2}+(\sqrt{1-b^{2}%
}-\sqrt{1-\tilde{b}^{2}})^{2},\end{split}
\end{aligned}\]
which is $1-\Theta(1/n^{3})$ if $b=\tilde{b}$, and at least $1+\Omega(1/n^{2})$ if $|b-\tilde{b}|\geq\Omega(1/n)$.
\par
\par
% END S2.SS0.SSS0.P0.SPx1.p2
% BEGIN S2.SS0.SSS0.P0.SPx1.p3
For item 3, observe that $\|q_{bc}-t_{\tilde{c}}\|^{2}_{2}=(\tilde{c}-c)^{2}+1,$ which is exactly 1 if $c=\tilde{c}$, and at least $1+\Omega(1/n^{10})$ if $|c-\tilde{c}|\geq\Omega(\varepsilon/n)$. Item 4 is easy to check. For item 5, observe that $\|s_{a}-t_{c}\|_{2}\geq 1+a\cos\delta+c\geq 1+\Omega(1/n^{4})$, whereas $\|s_{a}-q_{bc}\|_{2}$ and $\|p_{ab}-t_{c}\|_{2}$ are both $\sqrt{2}\pm o(1)$. $\hfill\blacktriangleleft$
\par
\par
% END S2.SS0.SSS0.P0.SPx1.p3
% BEGIN S2.SS0.SSS0.P0.SPx1.p4
In the real RAM model, balls $\mathcal{B}$ can clearly be constructed in $O(n+m)=\tilde{O}(n)$ time. In the RAM model, it suffices to approximate all coordinates up to $O(1/n^{10})$ additive error, which can be done in $\tilde{O}(1)$ time (we could actually pick “nice” choices of $b\in B$ so that $b$ and $\sqrt{1-b^{2}}$ are both rational, and also a nice $\delta$ so that $\cos\delta$ and $\sin\delta$ are both rational). Thus, the reduction can be done in $\tilde{O}(n)$ time. Theorem 9 follows from the following lemma.
\par
\par
% END S2.SS0.SSS0.P0.SPx1.p4
% BEGIN Thmtheorem12
\par\medskip\noindent\textbf{Lemma 12.}\par
Let $K$ be the intersection graph of the set of balls $\mathcal{B}$. Then $G$ has diameter at most $2$ if and only if $K$ has diameter at most $3$.
\par
\par
\par
% END Thmtheorem12
% END S2.SS0.SSS0.P0.SPx1
% BEGIN S2.SS0.SSS0.P0.SPx2
\par\medskip\noindent\textbf{Proof.}\par
% BEGIN S2.SS0.SSS0.P0.SPx2.p1
Observe by item 4 of Observation 11 that the balls in the same set $\mathcal{B}_{i}$ induce a clique in $K$. By item 5 of Observation 11, there is no edge between two balls in $\mathcal{B}_{i}$ and $\mathcal{B}_{j}$ when $|j-i|\geq 2$. Thus, $K$ has diameter at most $3$ if and only if for every ball $\llbracket s_{a}\rrbracket\in\mathcal{B}_{1}$ and every ball $\llbracket t_{c}\rrbracket\in\mathcal{B}_{4}$, $d_{K}(s_{a},t_{c})\leq 3$. Here we slightly abuse the notation by using the center of the ball to denote the corresponding vertex in $K$.
\par
\par
% END S2.SS0.SSS0.P0.SPx2.p1
% BEGIN S2.SS0.SSS0.P0.SPx2.p2
Suppose that $G$ has diameter at most $2$. Let $P=(a,b,c)$ be any path of length $2$ in $G$. Then by Observation 11, $(s_{a},p_{ab},q_{bc},t_{c})$ is a path of length 3 in $K$. Thus, $K$ has diameter at most 3 by the above observation. For the other direction, if $K$ has diameter at most $3$, then for any two vertices $s_{a}$ and $t_{c}$ (such that $\llbracket s_{a}\rrbracket\in\mathcal{B}_{1}$ and $\llbracket t_{c}\rrbracket\in\mathcal{B}_{4}$), a path of length 3 between $s_{a}$ and $t_{c}$ in $K$ must be of the form $(s_{a},p_{\tilde{a}b},q_{\tilde{b}c},t_{c})$. Then again by Observation 11, $a=\tilde{a}$ and $b=\tilde{b}$, and therefore, $(a,b,c)$ is a path of length $2$ in $G$. $\hfill\blacktriangleleft$
\par
\par
% END S2.SS0.SSS0.P0.SPx2.p2
% END S2.SS0.SSS0.P0.SPx2
% END S2
% BEGIN bib
\par\medskip\noindent\textbf{References}\par
\begin{enumerate}
\item[]  [1]  A. Karim Abu-Affash, Paz Carmi, Anil Maheshwari, Pat Morin, Michiel Smid, and Shakhar Smorodinsky.  Approximating maximum diameter-bounded subgraph in unit disk graphs.  Discret. Comput. Geom., 66(4):1401–1414, 2021.  doi:10.1007/s00454-021-00327-y.  
\item[]  [2]  Pankaj K Agarwal.  Range searching.  In Handbook of discrete and computational geometry, pages 1057–1092. Chapman and Hall/CRC, 2017.  doi:10.1201/9781315119601.  
\item[]  [3]  Sujoy Bhore, Timothy M. Chan, Zhengcheng Huang, Shakhar Smorodinsky, and Csaba D. Tóth.  Sparse bounded hop-spanners for geometric intersection graphs.  In Oswin Aichholzer and Haitao Wang, editors, 41st International Symposium on Computational Geometry, SoCG 2025, Kanazawa, Japan, June 23-27, 2025, volume 332 of LIPIcs, pages 17:1–17:15. Schloss Dagstuhl – Leibniz-Zentrum für Informatik, 2025.  doi:10.4230/LIPIcs.SOCG.2025.17.  
\item[]  [4]  Nicolas Bousquet and Stéphan Thomassé.  VC-dimension and Erdős–Pósa property.  Discret. Math., 338(12):2302–2317, 2015.  doi:10.1016/j.disc.2015.05.026.  
\item[]  [5]  Karl Bringmann, Sándor Kisfaludi-Bak, Marvin Künnemann, André Nusser, and Zahra Parsaeian.  Towards sub-quadratic diameter computation in geometric intersection graphs.  In Xavier Goaoc and Michael Kerber, editors, 38th International Symposium on Computational Geometry, SoCG 2022, Berlin, Germany, June 7-10, 2022, volume 224, pages 21:1–21:16. Schloss Dagstuhl – Leibniz-Zentrum für Informatik, 2022.  doi:10.4230/LIPIcs.SoCG.2022.21.  
\item[]  [6]  Sergio Cabello.  Subquadratic algorithms for the diameter and the sum of pairwise distances in planar graphs.  ACM Trans. Algorithms, 15(2):21:1–21:38, December 2019.  doi:10.1145/3218821.  
\item[]  [7]  Timothy M. Chan.  Finding triangles and other small subgraphs in geometric intersection graphs.  In Nikhil Bansal and Viswanath Nagarajan, editors, Proceedings of the 2023 ACM-SIAM Symposium on Discrete Algorithms, SODA 2023, Florence, Italy, January 22-25, 2023, pages 1777–1805. SIAM, 2023.  doi:10.1137/1.9781611977554.CH68.  
\item[]  [8]  Timothy M. Chan, Hsien-Chih Chang, Jie Gao, Sándor Kisfaludi-Bak, Hung Le, and Da Wei Zheng.  Truly subquadratic time algorithms for diameter and related problems in graphs of bounded VC-dimension.  In 66th Annual IEEE Symposium on Foundations of Computer Science (FOCS), pages 2728–2765, 2025.  doi:10.48550/arXiv.2510.16346.  
\item[]  [9]  Timothy M. Chan, Hsien-Chih Chang, Jie Gao, Sándor Kisfaludi-Bak, Hung Le, and Da Wei Zheng.  Charting the diameter computation landscape on intersection graphs in 2D, 2026.  Under preparation.  
\item[]  [10]  Timothy M. Chan, Qizheng He, and Yuancheng Yu.  On the fine-grained complexity of small-size geometric set cover and discrete k-center for small k.  In Kousha Etessami, Uriel Feige, and Gabriele Puppis, editors, 50th International Colloquium on Automata, Languages, and Programming, ICALP 2023, Paderborn, Germany, July 10-14, 2023, volume 261 of LIPIcs, pages 34:1–34:19. Schloss Dagstuhl – Leibniz-Zentrum für Informatik, 2023.  doi:10.4230/LIPIcs.ICALP.2023.34.  
\item[]  [11]  Hsien-Chih Chang, Jie Gao, and Hung Le.  Computing diameter+2 in truly-subquadratic time for unit-disk graphs.  In Wolfgang Mulzer and Jeff M. Phillips, editors, 40th International Symposium on Computational Geometry, SoCG 2024, Athens, Greece, June 11-14, 2024, volume 293, pages 38:1–38:14. Schloss Dagstuhl – Leibniz-Zentrum für Informatik, 2024.  Arxiv version https://arxiv.org/abs/2401.12881.  doi:10.4230/LIPIcs.SoCG.2024.38.  
\item[]  [12]  Victor Chepoi, Bertrand Estellon, and Yann Vaxès.  Covering planar graphs with a fixed number of balls.  Discret. Comput. Geom., 37(2):237–244, 2007.  doi:10.1007/s00454-006-1260-0.  
\item[]  [13]  Guillaume Ducoffe, Michel Habib, and Laurent Viennot.  Diameter, eccentricities and distance oracle computations on H-minor free graphs and graphs of bounded (distance) vapnik-chervonenkis dimension.  SIAM J. Comput., 51(5):1506–1534, 2022.  doi:10.1137/20M136551X.  
\item[]  [14]  Lech Duraj, Filip Konieczny, and Krzysztof Potępa.  Better diameter algorithms for bounded VC-dimension graphs and geometric intersection graphs.  In Timothy M. Chan, Johannes Fischer, John Iacono, and Grzegorz Herman, editors, 32nd Annual European Symposium on Algorithms, ESA 2024, Royal Holloway, London, United Kingdom, September 2-4, 2024, volume 308, pages 51:1–51:18. Schloss Dagstuhl – Leibniz-Zentrum für Informatik, 2024.  doi:10.4230/LIPIcs.ESA.2024.51.  
\item[]  [15]  Paweł Gawrychowski, Haim Kaplan, Shay Mozes, Micha Sharir, and Oren Weimann.  Voronoi diagrams on planar graphs, and computing the diameter in deterministic $\tilde{O}(n^{5/3})$ time.  SIAM J. Comput., 50(2):509–554, 2021.  doi:10.1137/18M1193402.  
\item[]  [16]  Hung Le and Christian Wulff-Nilsen.  VC set systems in minor-free (di)graphs and applications.  In David P. Woodruff, editor, Proceedings of the 2024 ACM-SIAM Symposium on Discrete Algorithms, SODA 2024, Alexandria, VA, USA, January 7-10, 2024, pages 5332–5360. SIAM, 2024.  doi:10.1137/1.9781611977912.192.  
\item[]  [17]  Jason Li and Merav Parter.  Planar diameter via metric compression.  In Moses Charikar and Edith Cohen, editors, Proceedings of the 51st Annual ACM SIGACT Symposium on Theory of Computing, STOC 2019, Phoenix, AZ, USA, June 23-26, 2019, pages 152–163, Phoenix, AZ, USA and New York, NY, USA, 2019. ACM.  doi:10.1145/3313276.3316358.  
\item[]  [18]  Jiří Matoušek.  Range searching with efficient hierarchical cuttings.  Discret. Comput. Geom., 10:157–182, 1993.  doi:10.1007/BF02573972.  
\item[]  [19]  Liam Roditty and Virginia Vassilevska Williams.  Fast approximation algorithms for the diameter and radius of sparse graphs.  In Dan Boneh, Tim Roughgarden, and Joan Feigenbaum, editors, Symposium on Theory of Computing Conference, STOC’13, Palo Alto, CA, USA, June 1-4, 2013, pages 515–524. ACM, 2013.  doi:10.1145/2488608.2488673.  
\item[]  [20]  Norbert Sauer.  On the density of families of sets.  J. Comb. Theory A, 13(1):145–147, 1972.  doi:10.1016/0097-3165(72)90019-2.  
\item[]  [21]  Saharon Shelah.  A combinatorial problem; stability and order for models and theories in infinitary languages.  Pacific Journal of Mathematics, 41(1):247–261, 1972.  doi:10.2140/pjm.1972.41.247.  
\end{enumerate}
% END bib
